% !PRINTED none
% Section-opening page: no page number is visible in the running head of any of
% the three bands. Page carries the Abschnitt title, the \S.\,63 sub-heading, the
% opening paragraphs of the section, and a footnote citing the literature.

\gesperrt{Ninth Section.}

\gesperrt{Congruence Groups.}

\S.\,63.

\gesperrt{Function Congruences.}

From the linear substitutions, whose formation and composition we have come to know in the earlier sections, yet another, wholly different kind of finite group can be derived: the \gesperrt{congruence groups}.
% [sic] Weber prints the plural verb "lassen" with the singular head "eine ... Art"; the English gives the intended grammatical sense.

These congruence groups have a manifold interest. They stem in the first place from the theory of elliptic functions, where they present themselves as \gesperrt{Galois} groups of the modular equations. But they are also of importance for general group theory, because they give us a means of forming whole sequences of \gesperrt{simple groups}. This is all the more remarkable since, as we have seen in the fourth section, simple groups, at least among the lower degrees, are very rare.

This theory of congruence groups is not only greatly generalized; its governing laws also emerge much more clearly when one relies on an extension of the concept of congruence originating with \gesperrt{Gauss}, which has since been developed and applied several times to problems of group theory, especially by \gesperrt{Galois}\footnote{\gesperrt{Galois}, ``Sur la théorie des nombres''. Bulletin des sciences math. de Ferussac. 1830. (Cf. the note to \S.\,151 of the first volume.) --- \gesperrt{Schönemann}, Grundzüge einer allgemeinen Theorie der höheren Congruenzen, deren Modulus eine reelle Primzahl ist. (Crelle's Journ. f. Mathematik, Bd. 31, 1846.) --- \gesperrt{Dedekind}, Abriss einer Theorie der höheren Congruenzen in Bezug auf einen reellen Primzahlmodulus. (Crelle's Journ. f. Mathematik, Bd. 54, 1856.)}.
